\documentclass[runningheads]{llncs}
\usepackage[T1]{fontenc}
\usepackage{graphicx}
\usepackage{float}
\usepackage{booktabs}
\usepackage{tabularx}
\usepackage{color}
\usepackage[colorlinks=true, linkcolor=black, citecolor=black, urlcolor=blue]{hyperref}
\renewcommand\UrlFont{\color{blue}\rmfamily}
\urlstyle{rm}
\usepackage{soul}
\begin{document}
\title{MSc Project 1: Time-Series-Enhanced Predictive Process Monitoring}
\subtitle{Phase Review 3}
\titlerunning{Time-Series-Enhanced PPM: Phase Review 3}
\author{Michał Durkalec \and
    Konstantinos Sakkas \and
    Roman Ležovič
}
\authorrunning{M. Durkalec et al.}
\institute{Supervised Process Intelligence (SPI) Lab,\\
    RWTH Aachen University, Aachen, Germany\\
    \email{office@pads.rwth-aachen.de}}
\maketitle
\begin{abstract}
    This document reviews Sprint~3 of the project, covering its objectives, technical
    implementation, and retrospective. Building on the time-series feature foundation
    established previously, this sprint focused on enriching and optimising the
    time-series feature set, adding dimensionality reduction and feature-pruning support,
    and implementing the evaluation, selection, and explainability machinery needed to
    compare baseline and time-series-enhanced models. We expanded the active-case-count
    feature into rolling-window statistics and vectorised their computation, wired
    configurable PCA and feature pruning into the pipeline, and added plateau-based prefix
    selection, model comparison, a feature-impact comparison tool, and permutation-based
    feature importance. The notebook suite was reworked to expose these capabilities through
    the same public API as the command-line pipeline. The empirical results and the
    baseline-versus-enhanced model comparison produced by this machinery are reserved for a
    dedicated evaluation document in the next phase.
    \keywords{Predictive Process Monitoring \and Process Mining \and Time-Series Features
        \and Feature Selection \and Model Evaluation \and Explainability}
\end{abstract}

\section{Introduction \& Objectives}

\subsection{Sprint Goals}
The goal of Sprint~3 was to make the pipeline ready to answer the project's central
question --- whether time-series information improves predictive process monitoring ---
by completing the feature and evaluation infrastructure required for a rigorous comparison.
Concretely, the sprint aimed to (i)~broaden the set of time-series features beyond a single
active-case count and improve their computational efficiency, (ii)~add mechanisms to manage
feature dimensionality and prune uninformative features, (iii)~implement the selection logic
that identifies the most suitable prefix length and model for each task, and (iv)~provide the
explainability and comparison tooling needed to contrast log-only baselines with
time-series-enhanced models. In line with the project's incremental approach, this sprint
delivers the \emph{capability} to run these analyses; the resulting metrics and the model
comparison itself are deferred to a separate evaluation document in the next phase.

\subsection{Team Organization}
Work continued to be tracked through GitHub Issues, with each track captured as a parent
issue and individual tasks as linked sub-issues and pull requests, keeping the \texttt{dev}
branch in a working state throughout the sprint. Michal coordinated the sprint and led the
feature, model, and prefix selection tracks together with reporting and visualisation.
Konstantinos implemented the time-series vectorisation, feature pruning, model-comparison, and
notebook-revamp pull requests. Roman implemented PCA-based dimensionality reduction and added further log based features.

\subsection{Role Allocation}
Initial role allocation was done according to table \ref{tab:roles}, but with this sprint being the last one most roles responsibilities were already fulfilled and issues were assigned to people dynamically and without constraints.

\begin{table}[H]
    \caption{Role assignments for Sprints 1--3.}\label{tab:roles}
    \centering
    \begin{tabular}{@{}llll@{}}
        \toprule
        \textbf{Role}                      & \textbf{Sprint 1}                                        & \textbf{Sprint 2} & \textbf{Sprint 3} \\
        \midrule
        Project Coordinator / Scrum Master & Michal                                                   & Konstantinos      & Roman             \\
        Technical Lead / Architecture      & Roman                                                    & Michal            & Konstantinos      \\
        Data \& Preprocessing Lead         & Roman, Konstantinos                                      & Konstantinos      & Roman             \\
        Modelling \& Experimentation Lead  & Michal                                                   & Roman             & Michal            \\
        Evaluation \& Visualization Lead   & Konstantinos                                             & Michal            & Konstantinos      \\
        Quality \& Tooling (Tests, CI)     & Konstantinos                                             & Roman             & Roman             \\
        Documentation \& Reporting         & \multicolumn{3}{c}{Michał, Roman, Konstantinos (shared)}                                         \\
        \bottomrule
    \end{tabular}
\end{table}

\section{Technical Implementation}
Sprint~3 extended the modular pipeline established in earlier sprints without restructuring it:
every addition slots into the existing strategy-based stages defined in
\texttt{data/schemas.py} and \texttt{data/types.py}. The following subsections describe the new
time-series features and their vectorisation, the dimensionality-reduction and feature-pruning
support, the model and prefix selection logic, the comparison and explainability tooling, and
the supporting configuration and notebook changes. \autoref{tab:steps} maps each work item to
its issues and pull requests and the modules it introduced or extended.

\subsection{Time-Series Feature Expansion and Vectorisation}
The time-series feature track was extended from a single active-case indicator to a richer set
of derived signals. The \texttt{ActiveCaseCountFeature} computes the number of cases active at
the timestamp of a prefix's last event, while \texttt{preprocess\_time\_series} builds a
continuous hourly grid spanning the full log and \texttt{extract\_time\_series\_features}
computes rolling-window statistics over the active-case series. These features are aligned to
case prefixes by timestamp through \texttt{align\_features\_to\_prefixes}, preserving the strict
temporal correspondence required to avoid leakage. The window and statistic computations, which
were originally implemented with per-window Python loops, were rewritten using vectorised array
operations, removing a significant bottleneck in feature extraction.

\subsection{Dimensionality Reduction and Feature Pruning}
Two complementary mechanisms for controlling the feature space were added and exposed through
configuration. Principal Component Analysis is now an optional, configurable pipeline stage
governed by a \texttt{pca\_config} block (\texttt{active} and \texttt{keep\_variability}),
allowing a target fraction of retained variance to be specified per run. Independently, the
feature-engineering configuration gained a \texttt{drop\_features} list that removes named
columns from the training and test matrices before model fitting; this provides a direct,
reproducible way to prune features that prior importance analysis flags as uninformative.
Together these support the ``best feature set'' selection workflow for both log-based and
time-series features at the configuration level, without code changes.

\subsection{Prefix and Model Selection Logic}
Hyperparameter tuning is performed per model via \texttt{RandomizedSearchCV} in
\texttt{search\_hyperparams}, driven by each model's \texttt{param\_grid} and the shared
\texttt{search} configuration (iterations, cross-validation folds, optional search sample size).
For prefix selection, \texttt{\_detect\_plateau\_prefix} walks prefix lengths in ascending order
and identifies the point at which the relative improvement of the tracked metric falls below a
configurable threshold (default 5\%), marking the prefix length beyond which additional events
add little predictive value. \texttt{compare\_models} aggregates the per-prefix metrics into a
structured comparison that ranks models and records each model's plateau prefix, providing the
basis for the ``best model'' and ``best prefix'' selection for each task.

\subsection{Comparison and Explainability Tooling}
To contrast log-only baselines against time-series-enhanced models, a dedicated command-line
tool (\texttt{evaluation/evaluate\_feature\_impact.py}) loads the evaluation artifacts of two
runs, combines them into a long-format table tagged by source (baseline versus advanced), and
produces faceted, per-model plots of each metric across prefix lengths. Explainability is
provided by \texttt{evaluation/feature importance.py}, which computes permutation importance
both globally and per prefix length and renders importance bar charts, per-prefix heatmaps, and
importance trajectories. These tools generate the inputs for the forthcoming comparison; the
analysis of their output is intentionally left to the next phase.

\subsection{Configuration and Notebook Support}
Supporting changes improved usability and reproducibility. The command-line pipeline now derives
a default output directory from the configuration file path, so runs are written to predictable,
config-specific locations. The notebook suite was reworked into a coherent set
(\texttt{00\_data\_loading}, \texttt{01\_exploratory data analysis},
\texttt{02\_feature\_analysis}, \texttt{03\_model\_evaluation}) that drives the pipeline through
the same public API as the command-line entry point, demonstrating feature extraction,
dimensionality analysis, and the loading of evaluation artifacts without bespoke wrappers.

\begin{table}[H]
    \centering
    \caption{Sprint~3 implementation steps.}
    \label{tab:steps}
    \begin{tabularx}{\textwidth}{lXX}
        \toprule
        \textbf{Step} & \textbf{Description}                      & \textbf{Issues / PRs} \\
        \midrule
        1             & Expand and vectorise time-series features & \#79, \#87, \#88      \\
        2             & PCA-based dimensionality reduction        & \#76, \#84            \\
        3             & Feature pruning via configuration         & \#92                  \\
        4             & Plateau-based prefix and model selection  & \#80, \#81, \#90      \\
        5             & Baseline-vs-enhanced comparison tool      & \#82, \#93            \\
        6             & Permutation feature importance            & \#73                  \\
        7             & Metric reporting and visualisation        & \#78, \#91            \\
        8             & Config-derived output paths               & \#95                  \\
        9             & Notebook revamp                           & \#98 (\#97)           \\
        \bottomrule
    \end{tabularx}
\end{table}

\section{Review \& Retrospective}

\subsection{Successes}
The sprint completed its planned scope. The pipeline now supports an enriched, efficiently
computed time-series feature set, configurable dimensionality reduction and feature pruning, and
a full selection-and-explainability layer, all integrated without disturbing the existing modular
architecture. \autoref{tab:sprint3_issues} lists the issues and pull requests closed during the
sprint.

\begin{table}[H]
    \centering
    \caption{Sprint~3 issues and pull requests.}
    \label{tab:sprint3_issues}
    \begin{tabularx}{0.85\textwidth}{lX}
        \toprule
        \textbf{Issue / PR} & \textbf{Description}                                            \\
        \midrule
        \#73                & Implement feature importance                                    \\
        \#74                & Choose best time-series features                                \\
        \#75                & Choose best log-based features                                  \\
        \#76                & Implement dimensionality reduction with PCA                     \\
        \#78                & Visualize reports with metrics                                  \\
        \#79                & Add more time-series features                                   \\
        \#80                & Choose best prefix length                                       \\
        \#81                & Choose best model for both tasks                                \\
        \#82                & Compare model performance with and without time-series features \\
        \#87                & Vectorize time series and windows                               \\
        \#88                & Vectorize time series and windows (PR)                          \\
        \#90                & Choose best prefix and model (PR)                               \\
        \#91                & Visualize reports with metrics (PR)                             \\
        \#92                & Feature pruning (PR)                                            \\
        \#93                & Model comparison (PR)                                           \\
        \#95                & Choose default output path based on config path                 \\
        \#98                & Revamped notebooks (PR, \#97)                                   \\
        \bottomrule
    \end{tabularx}
\end{table}

\subsection{Difficulties}
A recurring tension during the sprint was the separation between
delivering selection and comparison \emph{machinery} and producing the actual results: several
``choose best'' issues required the supporting evaluation code to be in place before any
selection could be made, so the concrete choices of feature set, prefix length, and model are
deferred to the evaluation phase. The exploratory data analysis notebook originally scoped as
issue~\#89 was closed as \emph{not planned} after its content was absorbed into the reworked
notebook suite.

\section{Moving Forward}

\subsection{Future Planning}
With the feature, selection, and explainability infrastructure complete, the next phase shifts
from implementation to analysis. The planned work is:
\begin{itemize}
    \item \textbf{Evaluation and model comparison.} Run the baseline and time-series-enhanced
          configurations end-to-end and use the comparison and feature-importance tooling to
          assess whether time-series features improve predictive performance, analysing results
          across prefix lengths for both the outcome and remaining-time tasks.
    \item \textbf{Results document.} Produce the dedicated evaluation report presenting these
          findings, including the comparative analysis and feature-importance discussion.
    \item \textbf{Finalisation.} Consolidate documentation and the final report, and close all
          remaining open items.
\end{itemize}

\begin{thebibliography}{8}
    \bibitem{bpic2017}
    van Dongen, B.F.: BPI Challenge 2017. Eindhoven University of Technology (2017). \url{https://research.tue.nl/en/datasets/bpi-challenge-2017}
\end{thebibliography}
\end{document}